\section{Causal interventions and identifiability boundaries}
\label{sec:method}

The mechanism suggests interventions, but an intervention that changes an
attention result is not automatically a new positional encoding.  This section
separates three levels: a causal score intervention, a functionally equivalent
phase edit for one query, and a rotary geometry that can be reused by future
queries through one physical GQA cache.

\subsection{Same-path singleton intervention}

For a selected layer--head--token coordinate, we first execute the custom
attention graph with $\epsilon=0$, then repeat it with
$s_{l,h,j}\leftarrow s_{l,h,j}+\epsilon$.  Define the global pair margin under
that intervention as
\begin{equation}
  m^{(l,h,j)}(\epsilon)
  =m\!\left(\operatorname{do}
  [s_{l,h,j}\leftarrow s_{l,h,j}+\epsilon]\right),
  \qquad
  \Delta m^{\mathrm{FD}}_{l,h,j}(\epsilon)
  =m^{(l,h,j)}(\epsilon)-m^{(l,h,j)}(0),
  \label{eq:finite-intervention}
\end{equation}
and the first-order prediction is $\epsilon U_{l,h,j}$ from
\cref{eq:directed-utility}.  Comparing with the same-path no-op is essential:
native and instrumented BF16 kernels can differ slightly even when they retain
the same decision.  Target coordinates are ranked by the absolute product of
positive grid-envelope RoPE suppression and the exact answer-margin derivative
$\partial m/\partial s$.  This oracle deliberately favors decisive coordinates,
so the experiment tests causal closure rather than deployable token selection.
Matched-random coordinates preserve layer, head, and class but not decisive-token
identity; they are numerical controls, not a fully matched magnitude or
signal-to-noise baseline.

\subsection{A phase edit for one query is a logit edit}

Fix one query at position $t$, candidate support $\mathcal C$, and Values
$\{v_j\}$.  Write $\tau_{tj}=j-t$.  For frequency plane $r$, let
$\delta_{j,r}$ be a phase correction in radians, not a token displacement.
The native and phase-edited scores are
\begin{equation}
  s_j=\frac{\kappa^2}{\sqrt d}\sum_r
  q_r^\top R(\tau_{tj}\omega_r)k_{j,r},
  \qquad
  s'_j=\frac{\kappa^2}{\sqrt d}\sum_r
  q_r^\top R(\tau_{tj}\omega_r+\delta_{j,r})k_{j,r},
  \qquad
  b_j=s'_j-s_j.
  \label{eq:phase-bias}
\end{equation}

\begin{proposition}[Single-query functional non-identifiability]
\label{prop:phase-bias}
Under fixed $\mathcal C$ and $\{v_j\}$, the phase-edited attention output is
exactly the output obtained by adding logit bias $b$:
\begin{equation}
  \sum_{j\in\mathcal C}\softmax(s')_jv_j
  =\sum_{j\in\mathcal C}\softmax(s+b)_jv_j.
  \label{eq:phase-bias-equivalence}
\end{equation}
At the level of attention probabilities, the effect of $b$ is observable only
modulo $b\mapsto b+c\mathbf 1$.  From the aggregated Value output alone, it may
be even less identifiable when the Values are not affinely independent.
\end{proposition}

The proposition is algebraic: $s'=s+b$ by definition.  It does not say that
all positional geometries are equivalent.  It says that a gain demonstrated
only by a query-specific, pairwise phase edit with unchanged support and Values
cannot by itself identify a reusable new positional encoding.

\subsection{When can one GQA cache realize a rotary repair?}

Fix layer $l$, KV group $g$, a prefix epoch with positions $j\le T$, and a set
of future query events $u=(h,t)$ with $h\in\mathcal H_g$ and $t>T$.  Let
$\delta_{u,j,r}$ be the desired correction, in radians, to the native relative
phase on frequency plane $r$.  We call it \emph{single-cache rotary-operator
realizable} if there exist candidate-independent Query rotations
$\mathsf Q_{u,r}\in\mathrm{SO}(2)$ and prefix-frozen, group-shared Key rotations
$\mathsf K_{g,j,r}\in\mathrm{SO}(2)$ such that
\begin{equation}
  \begin{aligned}
  q^\star_{u,r}&=\mathsf Q_{u,r}R(t\omega_r)q_{u,r},&
  k^\star_{g,j,r}&=\mathsf K_{g,j,r}R(j\omega_r)k_{g,j,r},\\
  \mathsf Q_{u,r}^{\top}\mathsf K_{g,j,r}
  &=R(\delta_{u,j,r})&&\quad\forall(u,j,r).
  \end{aligned}
  \label{eq:cache-operator}
\end{equation}
followed by a standard dot product with no pairwise phase operator.  The last
equality is the target constraint: it requires the two reusable operators to
realize the specified correction, rather than merely to exist.

\begin{theorem}[Single-cache rotary realizability]
\label{thm:cache-realizability}
The corrections are single-cache rotary-operator realizable if and only if
there exist $\beta_{g,j,r}$ and $\alpha_{u,r}$ such that
\begin{equation}
  \delta_{u,j,r}\equiv\beta_{g,j,r}-\alpha_{u,r}\pmod{2\pi}.
  \label{eq:separable-phase}
\end{equation}
Equivalently, for any two query events $u,v$ in the same KV group and any two
prefix positions $i,j$,
\begin{equation}
  \delta_{u,i,r}+\delta_{v,j,r}
  -\delta_{u,j,r}-\delta_{v,i,r}\equiv0\pmod{2\pi}.
  \label{eq:four-cycle}
\end{equation}
\end{theorem}

Every element of $\mathrm{SO}(2)$ is a rotation $R(\cdot)$, so
\cref{eq:cache-operator} and the group law yield $\delta=\beta-\alpha$.
Conversely, any such factorization constructs
$\mathsf Q_{u,r}=R(\alpha_{u,r})$ and
$\mathsf K_{g,j,r}=R(\beta_{g,j,r})$; separability is then equivalent to
vanishing four-cycles.
The full construction is in \cref{app:realizability}.  The factorization has a
common-phase gauge: adding the same phase to $\alpha$ and $\beta$ changes no
$\delta$.  The rectangular support is also important; for a sparse bipartite
support, the alternating signed phase sum must vanish on every graph cycle, and
each disconnected component has an independent common-phase gauge.
This theorem characterizes the phase operator, not accidental score equality on
a finite set of Q/K vectors, where zero amplitudes or cosine symmetries can add
degeneracies.  Because four Query heads share each Qwen3-8B KV head,
$\beta_{g,j,r}$ must be shared by those heads.  A violation requires per-query
Key work, multiple cached bases, or an explicit pairwise score/router operation.

\subsection{Diagnostic retrieval and a feasible design space}

We retain an exact pre-RoPE proposal followed by native post-RoPE sparse
attention as a diagnostic intervention in \cref{app:diagnostic-retrieval}.  It
tests whether semantically relevant keys suppressed at admission can still be
recovered, but it is neither the main contribution nor an end-to-end speedup:
the current implementation scans the full pre-RoPE history.  Closely related
proposal/consumption and local-memory structures already appear in SALS,
LongHeads, and InfLLM~\citep{mu2025sals,lu2024longheads,xiao2024infllm}.

Theorem~\ref{thm:cache-realizability} does expose a possible future design
space.  A cache-compatible semantic phase may take the factorized form
\begin{equation}
  \beta_{l,g,j,r}=f^K_{l,g,r}(x_{l,j}),
  \qquad
  \alpha_{l,h,t,r}=f^Q_{l,h,r}(x_{l,t}).
  \label{eq:semantic-phase-space}
\end{equation}
This is a theorem-permitted family, not a method trained or validated here.
It is also adjacent to content-dependent repositioning and head-adaptive RoPE
~\citep{li2026repo,wang2026adarope}.  A pair-dependent local/remote gate
generally breaks separability.  A static layer or KV-group partition can keep
exact local RoPE only in designated native layers or groups; it cannot make the
same affected head simultaneously pairwise-local-exact and remote-semantic.
That within-head requirement needs two cached bases or an explicit pairwise
router, with extra compute and possibly extra cache when multiple Key bases are
selected.
